\documentclass{article}
\usepackage[margin=1in, a4paper]{geometry}
\usepackage{amsmath, amssymb, amsfonts}
\usepackage{graphicx}
\usepackage{xcolor}
\usepackage{listings}
\usepackage{booktabs}
\usepackage[utf8]{inputenc}
\usepackage[T1]{fontenc}
\usepackage{hyperref}
\hypersetup{colorlinks=true, urlcolor=blue, linkcolor=blue}
\usepackage{enumitem}
\usepackage{float}

\definecolor{codegreen}{rgb}{0,0.6,0}
\definecolor{codegray}{rgb}{0.5,0.5,0.5}
\definecolor{codepurple}{rgb}{0.58,0,0.82}
\definecolor{backcolour}{rgb}{0.99,0.99,0.99}
% Professional color palette for academic publishing
\definecolor{myblue}{cmyk}{0.8,0.4,0,0.1}
\definecolor{mygreen}{cmyk}{0.7,0,0.5,0.2}
\definecolor{myorange}{cmyk}{0,0.5,0.8,0.1}
\definecolor{mygray}{cmyk}{0,0,0,0.4}
\definecolor{arrowcolor}{cmyk}{0.9,0.5,0,0.3}
\definecolor{nodecolor}{cmyk}{0.6,0.2,0,0.05}
\definecolor{framecolor}{cmyk}{0.4,0.1,0,0.1}

\lstdefinestyle{mystyle}{
    backgroundcolor=\color{backcolour},   
    commentstyle=\color{codegreen},
    keywordstyle=\color{magenta},
    numberstyle=\tiny\color{codegray},
    stringstyle=\color{codepurple},
    basicstyle=\ttfamily\footnotesize,
    breakatwhitespace=false,         
    breaklines=true,                 
    captionpos=b,                    
    keepspaces=true,                 
    numbers=left,                    
    numbersep=5pt,                  
    showspaces=false,                
    showstringspaces=false,
    showtabs=false,                  
    tabsize=2
}
\lstset{style=mystyle}

\title{The Newsvendor Problem}
\author{The Logistics \& Supply Chain Problem-Solving Playbook}
\date{}

\begin{document}
\maketitle

\section{The Newsvendor Problem}

The newsvendor problem represents one of the most fundamental challenges in inventory management and stochastic optimization. At its core, it addresses the eternal dilemma faced by decision-makers who must commit to inventory levels before demand uncertainty resolves itself. This problem extends far beyond the traditional newsstand scenario to encompass seasonal fashion merchandise, perishable goods, event-specific products, and any situation where replenishment is either impossible or prohibitively expensive during the selling period.

\subsection{The Scenario: Balancing Risk in the Face of Uncertain Demand}

Marina Rodriguez manages procurement for \textbf{Seasonal Sportswear Solutions}, a company specializing in high-end winter sports equipment. Every September, she faces a critical decision that will determine the company's profitability for the upcoming ski season: how many units of their premium ski jackets should she order from manufacturers in Southeast Asia?

The challenge is multifaceted and unforgiving. \textbf{Manufacturing lead times} stretch across four months, meaning orders must be placed by early October to arrive before the peak December-February selling season. Once the season begins, no additional inventory can be procured due to production scheduling constraints and shipping logistics. Each jacket costs \$180 to manufacture and ship, retails for \$320, and any unsold units must be liquidated at \$90 during the post-season clearance.

Historical data reveals that demand follows a complex pattern influenced by weather predictions, economic conditions, competitor actions, and fashion trends. Last season's demand ranged from 850 units (during an unusually warm winter) to 1,240 units (when a celebrity endorsement drove unexpected popularity). Marina knows that \textbf{stockouts} during peak season mean lost sales that cannot be recovered, while \textbf{overstock} situations tie up capital and erode margins through forced liquidation.

The newsvendor problem captures this essential tension between \textbf{underage costs} (the opportunity cost of not having enough inventory) and \textbf{overage costs} (the holding and liquidation costs of excess inventory). The optimal solution requires balancing these competing risks to maximize expected profit in the face of demand uncertainty.

This problem structure appears throughout supply chain management: pharmaceutical companies ordering flu vaccines before knowing the severity of the upcoming season, publishers determining print runs for books, farmers deciding crop allocations before knowing weather conditions, and manufacturers setting production schedules for seasonal products. The mathematical framework developed for the newsvendor problem provides the foundation for addressing these diverse but structurally similar challenges.

\subsection{Tier 1: The Pen \& Paper Method (Stochastic Programming Formulation)}

The newsvendor problem finds its mathematical foundation in stochastic programming, where we optimize decisions under uncertainty. We begin by establishing the formal mathematical framework that captures the essential trade-offs between underage and overage costs.

\textbf{Mathematical Formulation:}

Let us define the fundamental components of our stochastic programming model:

\textbf{Parameters:}
\begin{itemize}
\item $D$ = Random demand (with known probability distribution)
\item $p$ = Selling price per unit
\item $c$ = Cost per unit
\item $s$ = Salvage value per unsold unit (where $s < c < p$)
\item $f(d)$ = Probability density function of demand
\item $F(d)$ = Cumulative distribution function of demand
\end{itemize}

\textbf{Decision Variable:}
\begin{itemize}
\item $Q$ = Order quantity (inventory level to establish before demand realization)
\end{itemize}

\textbf{Stochastic Programming Objective Function:}

The expected profit function incorporates the stochastic nature of demand:

\begin{align}
\Pi(Q) &= \mathbb{E}[\text{Revenue}] - \mathbb{E}[\text{Costs}]\\
&= \mathbb{E}[p \cdot \min(D, Q) + s \cdot \max(0, Q-D)] - c \cdot Q
\end{align}

Expanding this expectation over the demand distribution:
\begin{align}
\Pi(Q) &= \int_0^Q p \cdot d \cdot f(d) \, dd + \int_Q^{\infty} p \cdot Q \cdot f(d) \, dd \\
&\quad + \int_0^Q s \cdot (Q-d) \cdot f(d) \, dd - c \cdot Q
\end{align}

\textbf{Critical Fractile Solution:}

Taking the derivative with respect to $Q$ and setting equal to zero:
\begin{align}
\frac{d\Pi(Q)}{dQ} &= p \cdot F(Q) + s \cdot F(Q) - s - c = 0\\
F(Q^*) &= \frac{p - c}{p - s} = \frac{c_u}{c_u + c_o}
\end{align}

where:
\begin{itemize}
\item $c_u = p - c$ (underage cost: lost margin per unit short)
\item $c_o = c - s$ (overage cost: net loss per unit overstocked)
\end{itemize}

The optimal order quantity $Q^*$ satisfies the \textbf{critical fractile condition}: we should order until the probability of selling the marginal unit equals the ratio of underage cost to total shortage-related costs.

\begin{figure}[htbp]
\centering
\includegraphics[width=\textwidth]{../figures-png/line55.fig1.png}
\caption{Critical Fractile Solution for the Newsvendor Problem}
\end{figure}

\textbf{Concrete Example:}

Consider Marina's ski jacket procurement decision with the following parameters:
\begin{itemize}
\item Selling price: $p = \$320$
\item Unit cost: $c = \$180$  
\item Salvage value: $s = \$90$
\item Demand follows Normal distribution: $D \sim N(1000, 150^2)$
\end{itemize}

Calculate costs:
\begin{itemize}
\item Underage cost: $c_u = \$320 - \$180 = \$140$
\item Overage cost: $c_o = \$180 - \$90 = \$90$
\end{itemize}

Critical fractile: $F(Q^*) = \frac{140}{140 + 90} = \frac{140}{230} = 0.6087$

For a normal distribution with $\mu = 1000$ and $\sigma = 150$:
$Q^* = 1000 + 150 \cdot \Phi^{-1}(0.6087) = 1000 + 150 \cdot (0.273) = 1041$ units

\textbf{Verification:}
With $Q^* = 1041$ units:
\begin{itemize}
\item Expected sales: $\mathbb{E}[\min(D, 1041)] = 998.7$ units
\item Expected leftover: $1041 - 998.7 = 42.3$ units
\item Expected profit: $\$320 \times 998.7 + \$90 \times 42.3 - \$180 \times 1041 = \$142,601$
\end{itemize}

\subsection{Tier 2: The Classic Heuristic (Priority-Based Marginal Analysis)}

While the critical fractile provides the theoretical optimum, practical implementation often requires heuristic approaches that can handle discrete demand distributions, computational constraints, or situations where the exact distribution is unknown. The priority-based marginal analysis approach evaluates each potential unit of inventory based on its expected contribution to profit.

\textbf{Algorithm Theory:}

The marginal analysis heuristic operates on a simple principle: order an additional unit if and only if the expected profit from that unit exceeds zero. This creates a natural priority ordering based on the probability of sale for each incremental unit.

For the $q$-th unit of inventory:
\begin{align}
\text{Expected Marginal Profit}(q) &= P(\text{Demand} \geq q) \times (p-c) \\
&\quad + P(\text{Demand} < q) \times (s-c)\\
&= [1-F(q-1)] \times (p-c) + F(q-1) \times (s-c)
\end{align}

The algorithm continues adding inventory units as long as the expected marginal profit remains positive.

\textbf{Detailed Pseudocode:}

\lstinputlisting[language=Python, caption=Priority-Based Marginal Analysis Algorithm]{.././codes-python/line55.py1.py}

\begin{figure}[htbp]
\centering
\includegraphics[width=\textwidth]{../figures-png/line55.fig2.png}
\caption{Marginal Profit Analysis for Inventory Units}
\end{figure}

\textbf{Computational Complexity:}
The marginal analysis approach has time complexity $O(Q_{max})$ where $Q_{max}$ is the maximum feasible order quantity. Space complexity is $O(Q_{max})$ for storing marginal profit calculations. This linear complexity makes it highly suitable for real-time applications and large-scale implementations.

\textbf{Concrete Example Output:}

Applying the marginal analysis to Marina's ski jacket problem with discrete demand scenarios:

\lstinputlisting[language=Python, caption=Concrete Example Implementation]{.././codes-python/line55.py2.py}

\textbf{Output Results:}
\begin{verbatim}
Optimal Order Quantity: 1041 units
Critical Unit Marginal Profit: $0.12
Demand 850: Sales=850, Leftover=191, Profit=$116,690
Demand 920: Sales=920, Leftover=121, Profit=$129,490  
Demand 1000: Sales=1000, Leftover=41, Profit=$139,690
Demand 1080: Sales=1041, Leftover=0, Profit=$145,740
Demand 1150: Sales=1041, Leftover=0, Profit=$145,740
Demand 1240: Sales=1041, Leftover=0, Profit=$145,740
Expected Total Profit: $138,435.50
\end{verbatim}

The marginal analysis confirms the theoretical result while providing intuitive insight into the decision process: we continue ordering units as long as their expected contribution to profit remains positive.

\subsection{Tier 3: The Advanced Algorithm (Ant Colony Optimization for Multi-Period Extension)}

The classical newsvendor problem assumes a single selling period, but real-world inventory decisions often involve multiple periods with evolving demand patterns, learning opportunities, and dynamic pricing. Ant Colony Optimization provides a powerful framework for solving the multi-period newsvendor problem with complex constraints and multiple objectives.

\textbf{Multi-Period Newsvendor Problem Formulation:}

Consider a planning horizon of $T$ periods where inventory decisions must account for:
\begin{itemize}
\item Carryover inventory from previous periods
\item Demand learning and forecast updating
\item Seasonal demand patterns  
\item Capacity constraints and quantity discounts
\item Multiple products with substitution possibilities
\end{itemize}

The state space becomes: $S_t = (I_{t-1}, D_{1:t-1}, \theta_t)$ where $I_{t-1}$ is ending inventory, $D_{1:t-1}$ represents historical demand realizations, and $\theta_t$ captures exogenous factors (seasonality, market conditions).

\textbf{ACO Algorithm Design:}

Ant Colony Optimization models the multi-period inventory problem as a graph where nodes represent inventory states and edges represent ordering decisions. Each ant constructs a complete solution (inventory policy across all periods) using pheromone-guided probabilistic decisions.

\textbf{Detailed Pseudocode:}

\lstinputlisting[language=Python, caption=Ant Colony Optimization for Multi-Period Newsvendor]{.././codes-python/line55.py3.py}

\begin{figure}[htbp]
\centering
\includegraphics[width=\textwidth]{../figures-png/line55.fig3.png}
\caption{Ant Colony Network for Multi-Period Inventory Optimization}
\end{figure}

\textbf{Concrete Example Output:}

Consider Marina's business expanding to a 5-period planning horizon with seasonal demand patterns:

\lstinputlisting[language=Python, caption=Multi-Period ACO Implementation Results]{.././codes-python/line55.py4.py}

\textbf{Example Output:}
\begin{verbatim}
Iteration 0: Best=245867.43, Avg=198543.21
Iteration 20: Best=287234.56, Avg=251876.32
Iteration 40: Best=298765.12, Avg=276543.87
Iteration 60: Best=312456.78, Avg=289765.43
Iteration 80: Best=318923.45, Avg=298654.21
Iteration 100: Best=323456.89, Avg=305432.67
Iteration 120: Best=325678.34, Avg=310987.54
Iteration 140: Best=326789.45, Avg=314521.76

Multi-Period Inventory Policy:
Period 1: Order 650 units
Period 2: Order 875 units  
Period 3: Order 425 units
Period 4: Order 180 units
Period 5: Order 85 units
Expected Total Profit: $326,789.45
\end{verbatim}

The ACO approach successfully handles the complexity of multi-period planning, discovering inventory policies that balance carrying costs across periods while adapting to seasonal demand patterns and capacity constraints.

\subsection{Tier 4: The AI/ML/RL Augmentation Method (Unsupervised Learning for Demand Pattern Discovery)}

Traditional newsvendor models assume known demand distributions, but real-world demand often exhibits complex, hidden patterns that challenge parametric assumptions. Unsupervised learning techniques can discover these latent demand structures, enabling more robust inventory decisions under distributional uncertainty.

\textbf{Unsupervised Learning Framework:}

The approach combines multiple unsupervised learning techniques to extract actionable insights from historical demand data:

1. \textbf{Clustering Analysis}: Identify distinct demand regimes (high/low seasons, trend shifts)
2. \textbf{Dimensionality Reduction}: Extract key factors driving demand variation  
3. \textbf{Anomaly Detection}: Identify outlier periods for robust decision-making
4. \textbf{Pattern Mining}: Discover sequential demand patterns and seasonality

\textbf{Algorithm Implementation:}

\lstinputlisting[language=Python, caption=Unsupervised Learning Enhanced Newsvendor]{.././codes-python/line55.py5.py}

\begin{figure}[htbp]
\centering
\includegraphics[width=\textwidth]{../figures-png/line55.fig4.png}
\caption{Clustering and PCA Analysis of Demand Patterns}
\end{figure}

\textbf{Concrete Example Output:}

Applying unsupervised learning to Marina's expanded dataset with 3 years of historical demand:

\lstinputlisting[language=Python, caption=Unsupervised Learning Analysis Results]{.././codes-python/line55.py6.py}

\textbf{Example Output:}
\begin{verbatim}
Discovered Demand Clusters:
Cluster 0: Expected Demand = 1050, Optimal Qty = 1087, Weight = 0.603
Cluster 1: Expected Demand = 1320, Optimal Qty = 1369, Weight = 0.244  
Cluster 2: Expected Demand = 780, Optimal Qty = 809, Weight = 0.153

Weighted Optimal Inventory: 1135 units

Pattern Discovery Results:
- Identified 3 distinct demand regimes
- 15 anomalous periods detected
- Top 3 PCA components explain 84.7% of variance

Feature Importance Rankings:
1. Rolling Mean Demand: 0.892
2. Lag-1 Demand: 0.845  
3. Seasonal Sine Component: 0.634
4. Base Demand: 0.621
5. Temperature: 0.387
\end{verbatim}

The unsupervised learning approach successfully identifies hidden demand structures, providing Marina with a more robust inventory strategy that adapts to different market conditions while accounting for seasonal patterns and anomalous events.

\subsection{Tier 10: Proactive Demand \& Market Shaping}

Moving beyond reactive inventory management, Tier 10 represents a paradigm shift toward \textbf{proactive demand creation and market influence}. Rather than merely responding to uncertain demand, companies actively shape customer behavior, create new market segments, and influence purchasing patterns to align with their operational capabilities and strategic objectives.

\textbf{The Demand Shaping Framework:}

Proactive demand shaping operates on multiple behavioral and economic levers:

1. \textbf{Dynamic Pricing and Revenue Management}: Real-time price adjustments to smooth demand patterns
2. \textbf{Behavioral Nudging}: Psychological techniques to influence customer choice timing and quantities  
3. \textbf{Product Bundling and Substitution Management}: Strategic offering design to shift demand across products
4. \textbf{Supply-Demand Orchestration}: Coordinated marketing campaigns timed with inventory availability
5. \textbf{Market Segmentation and Personalization}: Targeted demand creation for specific customer groups

\textbf{Mathematical Framework for Demand Shaping:}

The proactive newsvendor problem extends the classical formulation to include controllable demand parameters:

Let $D(\mathbf{p}, \mathbf{a}, \mathbf{m})$ represent demand as a function of:
\begin{itemize}
\item $\mathbf{p}$ = Dynamic pricing vector across time periods
\item $\mathbf{a}$ = Marketing action intensities (advertising spend, promotion levels)
\item $\mathbf{m}$ = Message framing and nudge mechanisms
\end{itemize}

The integrated optimization problem becomes:
\begin{align}
\max_{Q, \mathbf{p}, \mathbf{a}, \mathbf{m}} \quad &\mathbb{E}[\Pi(Q, D(\mathbf{p}, \mathbf{a}, \mathbf{m}))] - C_m(\mathbf{a}) - C_p(\mathbf{p})\\
\text{subject to} \quad &\mathbf{p} \in [\mathbf{p}_{min}, \mathbf{p}_{max}]\\
&\sum_t a_t \leq A_{budget}\\
&\text{Brand consistency constraints}\\
&\text{Regulatory compliance requirements}
\end{align}

where $C_m(\mathbf{a})$ represents marketing costs and $C_p(\mathbf{p})$ captures pricing-related costs (customer acquisition, brand perception impacts).

\begin{figure}[htbp]
\centering
\includegraphics[width=\textwidth]{../figures-png/line55.fig5.png}
\caption{Integrated Demand Orchestration and Market Shaping Framework}
\end{figure}

\textbf{Behavioral Economics Integration:}

The demand shaping approach leverages established behavioral economics principles:

\textbf{Loss Aversion and Scarcity Effects:}
\begin{align}
D_{scarcity}(t) &= D_{base}(t) \times [1 + \lambda_{scarcity} \times S(Q_{remaining}(t))]\\
S(Q_{remaining}) &= \frac{1}{1 + \exp(\alpha(Q_{remaining} - Q_{threshold}))}
\end{align}

where $S(Q_{remaining})$ represents the scarcity effect function that amplifies demand as inventory approaches depletion.

\textbf{Price Anchoring and Reference Point Effects:}
\begin{align}
D_{anchored}(p_t) &= D_{base} \times \left[\frac{p_{reference}}{p_t}\right]^{\beta_{anchor}} \times A(p_t, p_{reference})\\
A(p_t, p_{reference}) &= 1 + \gamma \times \max(0, p_{reference} - p_t)
\end{align}

\textbf{Social Proof and Network Effects:}
\begin{align}
D_{social}(t) &= D_{individual}(t) \times [1 + \delta \times N_{recent\_purchases}(t)]\\
N_{recent\_purchases}(t) &= \sum_{i=t-w}^{t-1} \text{purchases}_i \times \exp(-\phi(t-i))
\end{align}

\textbf{Concrete Example: Marina's Advanced Demand Shaping Strategy}

Marina implements a comprehensive demand orchestration system for her ski jacket business:

\textbf{Dynamic Pricing Algorithm:}
\lstinputlisting[language=Python, caption=Demand Shaping Implementation]{.././codes-python/line55.py7.py}

\textbf{Expected Output:}
\begin{verbatim}
Demand Shaping Strategy Results:
Optimal Inventory: 1456 units
Expected Profit: $398,742.33

Monthly Pricing Strategy:
Month 1: $295.50   Month 7: $340.25
Month 2: $285.75   Month 8: $355.00  
Month 3: $275.25   Month 9: $345.50
Month 4: $320.00   Month 10: $325.75
Month 5: $335.75   Month 11: $310.25
Month 6: $350.00   Month 12: $280.00

Monthly Marketing Investment:
Month 1: $3,200    Month 7: $6,500
Month 2: $4,100    Month 8: $5,800
Month 3: $5,750    Month 9: $4,200
Month 4: $7,200    Month 10: $3,500
Month 5: $6,800    Month 11: $4,800  
Month 6: $8,500    Month 12: $5,650

Key Strategic Insights:
- 32% profit improvement vs. traditional newsvendor approach
- Dynamic pricing captures 18% additional value through behavioral effects
- Coordinated marketing timing increases demand predictability by 24%
- Scarcity-driven urgency reduces end-of-season clearance by 45%
\end{verbatim}

\textbf{Advanced Behavioral Mechanisms:}

The Tier 10 approach incorporates sophisticated behavioral interventions:

1. \textbf{Pre-commitment and Reservation Systems}: Customers commit to future purchases at discounted rates, improving demand predictability
2. \textbf{Gamification and Loyalty Programs}: Points, badges, and tier systems encourage desired purchasing patterns
3. \textbf{Social Influence Campaigns}: User-generated content and influencer partnerships create organic demand amplification
4. \textbf{Personalized Scarcity Messaging}: AI-driven individual targeting of scarcity and urgency messages
5. \textbf{Cross-Category Demand Orchestration}: Coordinated promotions across complementary products to optimize overall portfolio profitability

Marina's advanced system demonstrates how modern retailers move beyond passive inventory management toward active market creation and customer behavior orchestration, achieving superior financial performance while enhancing customer satisfaction through better-aligned supply and demand.

## Practice Questions

\textbf{Tier 1 Questions (Mathematical Formulation):}

1. **Basic Critical Fractile Calculation**: A bookstore orders romance novels for Valentine's Day. Each novel costs \$8, sells for \$15, and unsold copies can be returned for \$3. If demand follows a normal distribution with mean 200 and standard deviation 40, calculate the optimal order quantity using the critical fractile method.

2. **Multi-Product Extension**: A seasonal retailer carries three products: winter coats (cost \$120, price \$200, salvage \$60), boots (cost \$80, price \$140, salvage \$40), and scarves (cost \$25, price \$45, salvage \$10). Demands are normally distributed with means (150, 200, 300) and standard deviations (30, 40, 60) respectively. Calculate the optimal order quantities and explain how correlation between demands would affect the solution.

3. **Service Level Constraint**: Reformulate Marina's ski jacket problem with an additional constraint requiring a 95\% service level (probability of not stocking out). How does this change the optimal solution compared to the pure profit maximization approach?

\textbf{Tier 2 Questions (Heuristic Implementation):}

4. **Discrete Demand Scenarios**: A farmer must decide how many acres to plant with a specialty crop. The crop costs \$500 per acre to plant and maintain, yields 1000 pounds per acre, and sells for \$2.50 per pound. Unsold crop can be sold as animal feed for \$0.80 per pound. Demand scenarios are: 180,000 lbs (prob 0.2), 220,000 lbs (prob 0.3), 260,000 lbs (prob 0.3), 300,000 lbs (prob 0.2). Use marginal analysis to determine optimal acreage.

5. **Computational Complexity Analysis**: Compare the computational requirements of the critical fractile method versus marginal analysis for problems with: (a) continuous normal demand, (b) discrete demand with 50 scenarios, (c) multi-period problems with 10 periods. Provide time and space complexity analysis.

\textbf{Tier 3 Questions (Metaheuristic Application):}

6. **Multi-Period ACO Implementation**: Extend the Ant Colony Optimization approach to handle a 6-month fashion retailer problem where unsold inventory from period $t$ can be carried to period $t+1$ at a holding cost of \$2 per unit per period. Demand forecasts decrease by 15\% each period for fashion items. Design the pheromone update mechanism and solution construction process.

7. **Hybrid Algorithm Design**: Design a hybrid algorithm combining Ant Colony Optimization with local search improvement for the newsvendor problem with quantity discounts: orders of 500+ units get 5\% discount, 1000+ units get 12\% discount, 1500+ units get 18\% discount. Explain how the algorithm handles the discrete breakpoints in the cost structure.

\textbf{Tier 4 Questions (AI/ML Integration):}

8. **Reinforcement Learning Formulation**: Formulate the multi-period newsvendor problem as a Markov Decision Process. Define the state space (including inventory levels and demand history), action space (order quantities), and reward function. Explain how Q-learning could be applied to learn optimal policies without knowing the demand distribution.

9. **Ensemble Demand Forecasting**: Design an ensemble machine learning approach that combines: (a) time series forecasting (ARIMA), (b) neural network predictions, (c) clustering-based regime identification, and (d) external factor regression. Explain how to weight the different models and integrate uncertainty quantification for robust newsvendor decisions.

\textbf{Tier 10 Questions (Proactive Demand Shaping):}

10. **Behavioral Intervention Design**: A luxury hotel must decide room allocation for New Year's Eve. Design a comprehensive demand shaping strategy incorporating: dynamic pricing (prices can vary from \$800-\$2000), targeted marketing campaigns, scarcity messaging, and social proof mechanisms. The hotel has 200 rooms, expects base demand of 180 rooms with high uncertainty. Formulate as an integrated optimization problem.

11. **Market Creation Strategy**: A tech startup launches a new wearable device with unknown market demand. Design a proactive demand creation strategy that includes: influencer partnerships, pre-order campaigns with early-bird pricing, limited edition releases, and viral marketing mechanisms. How would you model the demand evolution as marketing actions take effect, and optimize the production quantity given a 6-month lead time?

12. **Cross-Category Orchestration**: A grocery chain wants to optimize produce ordering while using complementary products to shape demand. Fresh salmon (cost \$8/lb, price \$15/lb, salvage \$3/lb) has uncertain demand, but can be bundled with rice, vegetables, and sauces to create meal kits. Design the integrated optimization framework that simultaneously determines salmon orders and promotional strategies for complementary items to maximize total category profit.

\end{document}
